\chapter{A comparative test}
\label{ch:evaluation-detail}

Give an author a smoother paragraph and a reader may still ask the same
question. That is why ``the prose looked better'' is not an estimand, and a
lower detector score is not one either. The outcome we can test is the
attention and meaning cost of getting a technical document to a named reader's
decision. We write that outcome down, choose comparisons that can identify it,
and show how a small feasibility sample can teach us without
pretending to be a definitive efficacy trial.

\section{The target population and unit of analysis}
\label{sec:eval-population}

The pilot needs documents whose authors and readers know enough of the field to
evaluate an argument, but not so much project history that private labels become
self-explanatory. Candidate documents include research proposals, model notes,
investment memoranda, and technical survey chapters. Each one must have a named
decision, at least one meaning-bearing object, and a reader who can describe the
decision context.

The unit of treatment is a document revision path. The unit of reader outcome
is a reader--version encounter. The unit of protected-object safety is an
object--version comparison. Keeping these units distinct prevents a document
with many equations from receiving more weight merely because it has more
tokens, and prevents one reader who spends a long time on a single draft from
being mistaken for many independent observations.

Let $d$ index documents, $v\in\{0,1\}$ denote the incumbent and Humanvoice
paths, and $j$ index readers. The protocol fixes $K$ review cycles and $H$
calendar days before assignment. Define
\begin{equation}
  Y_{djv} = \left(T^H_{djv},\ A^H_{dv},\ Q_{djv},\ R^K_{dv},\
  D^H_{dv},\ M_{dv}\right),
  \label{eq:eval-outcome-vector}
\end{equation}
where $T^H$ is cumulative reader time through the horizon, $A^H$ is acceptance
by the horizon, $Q$ is reconstruction-task completion, $R^K$ is restricted
feedback rounds, $D^H$ is the terminal disposition, and $M$ counts unresolved
protected-object changes. The two primary contrasts are
\begin{equation}
  \tau_A=\mathbb E[A^H_{d1}-A^H_{d0}],
  \qquad
  \tau_T=\mathbb E[T^H_{dj1}-T^H_{dj0}].
  \label{eq:eval-primary-estimand}
\end{equation}

All assigned documents remain in those contrasts. Accepted, still in revision,
declined, author-withdrawn, and operationally incomplete documents appear as
separate dispositions. The analysis does not condition the primary estimate on
meaning safety after assignment, because that would remove precisely the
failures the product is meant to prevent. It reports $M$ as a safety endpoint
that can stop adoption regardless of favorable time or acceptance.

For the declared reader task, let $Q_{djv}=1$ when the reader identifies the
object, mechanism, evidence, qualification, and requested action within the
time budget. The contrast
\begin{equation}
  \tau_Q=\mathbb E[Q_{dj1}-Q_{dj0}]
  \label{eq:eval-completion-estimand}
\end{equation}
is reported with a confidence interval and the complete answer distribution.
The project should not trade a small time saving for a lower probability of a
correct reconstruction or a higher decline rate.

\section{The comparisons}
\label{sec:eval-arms}

Three arms separate the value of diagnosis from the value of the full workflow:

\begin{description}[style=nextline,leftmargin=37mm,labelindent=0mm]
\item[Incumbent] The author's ordinary revision process, with the usual
grammar, version-control, and expert-review tools.
\item[Diagnostic path] Humanvoice findings and protected comparison are
available, but the author receives no model-generated rewrite.
\item[Full path] The diagnostic path plus optional local suggestions, with the
same author control, protected comparison, and reader task.
\end{description}

The diagnostic arm is important. If the full path improves the result, the
project needs to know whether the improvement came from located attention,
from wording suggestions, or from extra review time induced by the interface.
If the diagnostic arm performs as well as the full path, a smaller and easier
to govern product may be the better investment. If both arms fail against the
incumbent, the product claim should be revised rather than rescued by a
surface metric.

The first feasibility run can use a within-document crossover when the author
can produce two protected versions and the reader can be blinded to the path.
For a stronger causal comparison, later documents should be randomized to an
arm and readers should be randomized to version order. A crossover is efficient
for learning but vulnerable to carryover: a reader who saw one version knows
the decision before seeing the other. The protocol records order and analyzes
it rather than assuming it away.

Figure~\ref{fig:paired-design} illustrates the counterbalanced form. The exact
allocation will depend on document arrival and reader availability, but the
same outcomes must be recorded in both orders.

\hvFigurePairedDesign

\begin{table}[H]
\centering
\small
\begin{tabularx}{\textwidth}{@{}p{0.21\textwidth}L p{0.27\textwidth}@{}}
\toprule
Design choice & What it identifies & Main threat \\
\midrule
Within-document crossover & Difference for the same source and author & Reader
carryover and learning \\
Document-level randomization & Average effect across independent documents &
Small sample and heterogeneous genres \\
Reader order randomization & Position-adjusted comparison of versions & Incomplete
blinding if the interface leaves traces \\
Diagnostic versus full path & Contribution of model suggestions beyond findings &
Different author effort and novelty effects \\
\bottomrule
\end{tabularx}
\caption{The feasibility design and the stronger comparative design answer
different questions.}
\label{tab:eval-design-choices}
\end{table}

\section{The reader instrument}
\label{sec:eval-reader-instrument}

The reader instrument begins with context, not a style rubric. A reader is told
the decision, the role they occupy, the time budget, and the information they
would normally have. The instrument then asks four forced-choice questions and
one open question:

\begin{enumerate}
\item Which action is the author requesting?
\item What mechanism connects the proposed change to the outcome?
\item Which observation or calculation is the main support?
\item What condition or qualification limits the inference?
\item What would you ask the author before acting?
\end{enumerate}

The first four answers are coded for correctness against a document-specific
answer key prepared before randomization. The open question is retained for
diagnosis and is independently classified as a request for definition,
evidence, mechanism, scope, or action. The answer key is not a hidden test of
whether the reader shares the author's conclusion. A reader can disagree with
the recommendation and still understand it.

Time is measured from the first page shown to the final answer, with pauses and
clarification requests recorded separately. A reader who stops because the
document is clear has a different experience from one who stops after giving
up. The protocol therefore asks for a short reason whenever the reader ends
early, the page or section at which the stop occurred, and a short quoted span
that triggered the question. It also records the missing relation the reader
was trying to recover and whether the reader consulted a reference or source link;
that behavior is part of the document's burden, not an error to be removed.

\section{Sample size and uncertainty}
\label{sec:eval-size}

The feasibility study is for design learning, so its sample size is chosen to
estimate variance and failure modes rather than to promise a conventional
five-percent result. Suppose the primary outcome is paired reader time, with
within-document difference $D_d=T_{d1}-T_{d0}$. An approximate standard error
for the mean difference is
\begin{equation}
  se(\bar D)=\frac{s_D}{\sqrt{n}},
  \label{eq:eval-se}
\end{equation}
where $s_D$ is the observed standard deviation and $n$ is the number of
documents with usable paired readings. A planning interval of half-width $h$
requires approximately
\begin{equation}
  n \approx \left(\frac{z_{1-\alpha/2}s_D}{h}\right)^2.
  \label{eq:eval-n}
\end{equation}
The pilot does not know $s_D$ in advance. The one live feasibility case cannot
estimate it. Baseline records and a small observational set can provide a rough
planning range; the feasibility run then tests the timing instrument. Seeded
safety fixtures expose several meaning failures without asking a live author to
create them. Those inputs produce a range of sample sizes for the comparative
study rather than a spurious single number.

For task completion, a paired binary difference is less stable than a time
difference. Report the full two-by-two table of reader outcomes and use a
paired interval or a randomization distribution rather than a normal
approximation when the sample is small. Heterogeneity is not a nuisance to be
averaged away. A workflow that helps a junior author and burdens a senior
author has a different business case from one with a uniform effect. Report
document genre, author experience, reader role, and assistance condition with
each estimate.

The feasibility report should include a precision table before data collection:

\begin{table}[H]
\centering
\small
\begin{tabularx}{\textwidth}{@{}p{0.24\textwidth}p{0.19\textwidth}L@{}}
\toprule
Quantity & Planning value & Why it is recorded \\
\midrule
Documents & At least one complete document in each selected genre slice &
Prevents a single friendly case from defining the product \\
Readers per version & Two or more independent readers where feasible &
Separates reader disagreement from document effect \\
Protected objects & A mixture of equations, numbers, citations, tables, and
qualifications & Tests the safety boundary rather than prose alone \\
Time budget & Declared before showing the document & Makes completion comparable
across versions \\
\bottomrule
\end{tabularx}
\caption{The feasibility sample is designed to expose variance and failure
modes before a larger effect study is priced.}
\label{tab:eval-precision-plan}
\end{table}

\section{Agreement and adjudication}
\label{sec:eval-agreement}

Reader judgments are not perfectly reproducible, and a product should not hide
that fact behind a majority vote. For each question, retain the individual
answer, the answer key, the reader's confidence, and any explanation. A second
reader independently codes the open question. Disagreement is then adjudicated
by an evaluation reviewer who did not write the diagnostic or edit the tested
version. A domain reviewer may resolve a dispute about technical meaning, but
does not see the workflow assignment until the coding is frozen.

For finding labels, report agreement by category and prevalence, not just one
overall percentage. A rare but consequential meaning failure can disappear in
an aggregate agreement statistic. For model judges, calibrate against the
human-labeled pairs and measure position flips, length sensitivity, and
self-preference. Zheng et al., Panickssery et al., and Dubois et al. show why
these checks are necessary \citep{zheng2023judging,panickssery2024llm,
dubois2024length}.

The adjudicator does not rewrite the document during coding. That separation
keeps the outcome independent of the proposed repair. If a reader asks for a
change that the answer key did not anticipate, the request is preserved as a
new qualitative category and the key is updated only for a subsequent round.
This is how the instrument learns without moving its target after seeing the
result.

\section{Threats to interpretation}
\label{sec:eval-threats}

Several threats are predictable.

\begin{description}[style=nextline,leftmargin=37mm,labelindent=0mm]
\item[Novelty] A reader may spend more time on a new interface simply because
it is unfamiliar. The protocol gives the incumbent arm an equivalent context
and records the first-session effect.
\item[Selection] Authors may submit documents that they believe are especially
well suited to the tool. The feasibility report describes the intake process
and keeps rejected documents in an eligibility log.
\item[Contamination] A reader may have seen an earlier version or the source
history. The reader form asks about prior exposure. The assigned encounter
remains in the main analysis with a contamination flag; a pre-specified
sensitivity analysis reports uncontaminated encounters separately.
\item[Learning] In a crossover, the second reading benefits from the first.
Order is randomized, and the primary feasibility interpretation treats the
paired difference as descriptive when carryover is substantial.
\item[Hawthorne effect] Authors may work more carefully because they know a
study is observing them. The incumbent arm receives the same timing and
documentation burden where possible.
\item[Domain mismatch] A result in an economics note may not transfer to a
mathematical appendix or an investment memorandum. Genre and reader role are
reported rather than pooled silently.
\end{description}

The jagged-frontier evidence suggests an additional interaction: assistance
may be helpful for a task the model can perform and harmful when the task
requires knowledge it lacks \citep{dellacqua2026jagged}. The protocol therefore
asks the author to identify which parts of the document are ordinary language
editing and which parts require domain reasoning. A positive result confined
to the first category is still useful, but it supports a narrower product.

\section{From observations to an economic decision}
\label{sec:eval-economic-decision}

The economic model in Equation~\eqref{eq:value} needs observations from the
comparative test: restricted reader time, author and operator time, the
distribution of terminal decisions, protected-object incidents, document
volume, and local loaded rates. Table
\ref{tab:value-sensitivity} converts them into break-even questions without
inventing a dollar forecast.

The feasibility build also has learning value, but the proposal does not assign
that value a decorative probability. The observable decision is simpler: did
the build reveal enough about feasibility, burden, and variance to make the
comparative study cheaper to price than before? The result memorandum lists
which uncertainties narrowed and which remained untouched. That account lets
the sponsor value the option using its own alternatives and budget.

\section{Pre-specified decision rules}
\label{sec:eval-rules}

At the end of the feasibility run, the sponsor receives four results together:
reader time and task completion, protected-object safety, author burden, and
operating cost. The recommendation follows a simple set of conditions:

\begin{enumerate}
\item \textbf{Proceed to a comparative study} when the complete path is
      usable, no unresolved protected-object change remains, the reader outcome
      is measurable, and the precision calculation supports a priced next
      sample.
\item \textbf{Revise the product} when the reader outcome is promising but a
      named parser, finding family, privacy assumption, or interface causes
      avoidable burden.
\item \textbf{Stop or redirect} when meaning correspondence cannot be made
      reliable, the reader cannot recover the decision, or the operating cost
      exceeds any plausible attention saving in the chosen beachhead.
\end{enumerate}

These rules are intentionally qualitative at the feasibility stage. The study
should not manufacture a numerical pass threshold before the baseline and
variance are known. Once the comparative sample is specified, the sponsor can
set a minimum practically important time saving and a maximum tolerable
meaning-error rate. The current document records that future decision rather
than filling it with invented precision.

\section{Reporting the result}
\label{sec:eval-reporting}

Write the result like a small empirical paper. State the population, eligibility,
arms, reader context, order randomization, missing observations, protected-
object rules, and analysis before showing estimates. Show individual document
trajectories and reader questions alongside means. Keep a source result, a
local calculation, and a project hypothesis distinct in ordinary prose.

The forms can sit in the appendix. The main result should answer three
questions in sequence: did the workflow change the reader's task, did it keep
the document's meaning intact, and was the change worth the author's, reader's,
and operator's time? A model score may support those answers, never replace
them. That order keeps the proposal a decision document rather than a
collection of attractive diagnostics.
